\documentclass[UTF8,12pt,a4paper]{ctexart}
\usepackage{geometry}
\usepackage{booktabs}
\usepackage{longtable}
\usepackage{threeparttable}
\usepackage{array}
\usepackage{amsmath}
\usepackage{graphicx}
\usepackage{pdflscape}
\usepackage{setspace}
\usepackage{hyperref}
\geometry{left=2.6cm,right=2.6cm,top=2.5cm,bottom=2.5cm}
\setstretch{1.35}
\hypersetup{colorlinks=true,linkcolor=black,citecolor=black,urlcolor=blue}
\newcommand{\sym}[1]{\ifmmode^{#1}\else\(^{#1}\)\fi}

\title{审计师行业专长与企业漂绿\\阶段性研究设计与基准回归报告}
\author{}
\date{\today}

\begin{document}
\maketitle

\begin{abstract}
本报告基于2021--2024年公司年度面板，阶段性考察会计师事务所行业专长与企业漂绿程度之间的关系。研究以连续漂绿指标为被解释变量，分别以事务所在特定行业的客户营业收入市场份额（IMS）及其10\%门槛虚拟变量（IMS\_D）衡量审计师行业专长。为保证逐步回归中系数变化不受样本变化干扰，所有模型统一使用8,129个完整观测，并依次加入财务、治理和审计控制变量，最后控制年度与行业固定效应。结果显示，IMS在各列均显著为负；IMS\_D在各列也显著为负。最终采用公司聚类标准误的双向固定效应模型中，IMS和IMS\_D的系数分别为$-0.8739$和$-0.1129$，分别在1\%和5\%水平显著。现阶段结果表明行业专长与企业漂绿程度存在稳定的负向条件相关关系，但尚未完成替代变量、缩尾、二元模型和内生性处理，因而不能作因果解释。
\end{abstract}

\section{研究背景与阶段性目标}

行业专长审计师通常拥有更丰富的行业知识和经验，其专业胜任能力和职业判断有助于发现客户潜在错报、提高财务信息透明度，并缓解信息不对称（DeFond and Francis, 2000；刘文军等，2010；Reichelt and Wang, 2010）。已有研究还发现，审计师行业专长与知情交易概率、资本市场定价效率以及企业投资决策相关（陈小林等，2013；王生年等，2018；李姝等，2021）。参照李姝等（2021）从解释变量、控制变量和固定效应逐层构建多元回归模型的研究设计，本阶段仅检验会计师事务所层面的行业专长是否与企业漂绿程度存在稳定关系。

本报告不扩展新的作用机制或因果识别内容。现阶段目标包括：第一，验证简单相关关系在加入可观测公司特征后是否仍然存在；第二，通过连续变量IMS与虚拟变量IMS\_D两种口径检验结果是否一致；第三，比较公司聚类标准误与HC3异方差稳健标准误下的统计推断；第四，识别当前结果中值得关注的问题，为后续稳健性检验提供依据。

\section{研究设计}

\subsection{样本与数据处理}

研究面板包含2021--2024年观测。行业专长以完整A股审计市场为计算母样本，制造业采用证监会二位行业代码，非制造业采用一位行业代码；年度--行业中只有一家有效公司的组合已剔除。营业收入增长率先在2020--2024年完整财务数据中计算，再按证券代码和年度匹配至研究面板，从而避免在漂绿子样本内部计算滞后收入造成的大量缺失。

原面板共有9,239个观测。连续漂绿指标、行业专长、所有控制变量、公司代码、年度及行业代码均有效的共同样本为8,129个公司年度观测，涉及2,669家公司、4个年度和47个行业专长分类。所有逐步回归均固定使用这8,129个观测，因而不同列核心系数的变化来自模型设定变化，而非样本构成变化。

\subsection{变量定义与构造}

\subsubsection{被解释变量}

企业漂绿程度以\textbf{Greenwashing连续\_基准}衡量。该指标由公司在所属行业年度内标准化后的绿色披露指标与环境绩效指标之差构造：
\begin{equation}
Greenwashing_{it}=GI\_z_{it}-ESGE\_z_{it}.
\end{equation}
其中，$GI\_z$为绿色披露指标的行业年度标准化值，$ESGE\_z$为环境绩效得分的行业年度标准化值。指标越大，表示企业披露表现相对于实际环境绩效越高，即漂绿程度越高。本阶段仅以连续漂绿指标进行OLS基准回归；面板中的二元变量Greenwashing尚未用于本阶段模型。

\subsubsection{核心解释变量}

借鉴市场份额法衡量审计师行业专长的思路（Krishnan, 2003；陈小林等，2013；王生年等，2018），本研究在会计师事务所层面构造IMS：
\begin{equation}
IMS_{akt}=\frac{\sum_{i\in(a,k,t)}Revenue_{it}}
{\sum_{i\in(k,t)}Revenue_{it}},
\end{equation}
其中，$a$表示会计师事务所，$k$表示行业，$t$表示年度；分子是事务所$a$在行业$k$审计客户的营业收入合计，分母是该行业全部有效上市公司的营业收入合计。IMS越高，说明该事务所在相应行业审计市场中的客户收入份额越高。

IMS\_D为行业专长虚拟变量：
\begin{equation}
IMS\_D_{akt}=\begin{cases}
1,& IMS_{akt}\geq 10\%,\\
0,& IMS_{akt}<10\%.
\end{cases}
\end{equation}
固定10\%门槛与王生年等（2018）采用的事务所行业专长定义一致。IMS和IMS\_D分别进入模型，不在同一回归中同时使用，以避免二者机械相关影响估计。

\subsubsection{控制变量}

参照相关研究控制公司基本特征、财务状况、治理结构和审计特征（王生年等，2018；李姝等，2021）。各变量的具体定义见表\ref{tab:variables}。

\begin{longtable}{p{2.7cm}p{2.4cm}p{8.4cm}}
\caption{变量定义}\label{tab:variables}\\
\toprule
变量名称 & 变量符号 & 定义与构造方式\\
\midrule
\endfirsthead
\toprule
变量名称 & 变量符号 & 定义与构造方式\\
\midrule
\endhead
企业漂绿程度 & Greenwashing & 行业年度标准化绿色披露指标$GI\_z$减去行业年度标准化环境绩效指标$ESGE\_z$；数值越大表示漂绿程度越高。\\
行业市场份额 & IMS & 某事务所在某年度某行业的客户营业收入合计，除以该年度该行业全部有效上市公司的营业收入合计。\\
行业专家 & IMS\_D & IMS大于或等于10\%取1，否则取0。\\
公司规模 & Size & 期末总资产的自然对数，即$\ln(TotalAssets)$。\\
资产负债率 & Lev & 期末总负债除以期末总资产。\\
总资产收益率 & ROA & 当期净利润除以期末总资产。\\
经营现金流 & CFO & 经营活动产生的现金流量净额除以期末总资产。\\
营业收入增长率 & Growth & 当期营业收入除以上一期营业收入再减1；上一期收入来自完整财务数据，且要求年度连续、上一期收入为正。\\
公司年龄 & FirmAgeLog & 公司年龄加1后取自然对数，即$\ln(1+FirmAge)$。\\
产权性质 & SOE & 国有企业取1，非国有企业取0。\\
两职合一 & Dual & 董事长与总经理由同一人担任取1，否则取0。\\
董事会规模 & BoardSize & 董事会人数的自然对数，即$\ln(DirectorNumber)$。\\
独立董事比例 & IndRatio & 独立董事人数除以董事会人数。\\
国际四大 & Big4 & 境内审计事务所属于国际四大取1，否则取0。\\
审计收费 & LnAuditFee & 年度审计费用的自然对数。\\
审计任期 & AuditorTenure & 当前审计事务所连续为公司提供审计服务的年数。\\
事务所变更 & IsChangeFirm & 当年发生会计师事务所变更取1，否则取0。\\
年度固定效应 & Year FE & 控制2021--2024年共同宏观年度冲击。\\
行业固定效应 & Industry FE & 按行业专长计算口径控制行业不随观测个体变化的平均差异。\\
\bottomrule
\end{longtable}

\subsection{模型设定}

参照李姝等（2021）在基准模型中加入控制变量、年度和行业固定效应的做法，本研究设定：
\begin{equation}
Greenwashing_{it}=\alpha_0+\alpha_1Expertise_{it}
+\sum_j\beta_j Controls_{jit}+Year_t+Industry_k+\varepsilon_{it},
\label{eq:baseline}
\end{equation}
其中，$Expertise$分别取IMS和IMS\_D，$Controls$表示前述控制变量。由于同一公司跨年度误差可能相关，主要统计推断采用公司层面聚类标准误；HC3异方差稳健标准误作为同一模型系数下的推断对照。

\section{描述性统计}

表\ref{tab:desc}报告8,129个共同样本的描述性统计。连续漂绿指标均值为0.0019、中位数为$-0.0547$，标准差为1.3310，显示公司之间漂绿程度差异较大。IMS均值为0.0662，中位数为0.0402；IMS\_D均值为0.2292，表示约22.92\%的公司年度观测由市场份额达到10\%门槛的行业专家事务所审计。

\begin{landscape}
\begin{longtable}{lrrrrrrrr}
\caption{共同样本变量描述性统计}\label{tab:desc}\\
\toprule
变量 & N & 均值 & 标准差 & 最小值 & 25\% & 中位数 & 75\% & 最大值\\
\midrule
\endfirsthead
\toprule
变量 & N & 均值 & 标准差 & 最小值 & 25\% & 中位数 & 75\% & 最大值\\
\midrule
\endhead
Greenwashing & 8,129 & 0.0019 & 1.3310 & -6.5004 & -0.8092 & -0.0547 & 0.9019 & 5.0259\\
IMS & 8,129 & 0.0662 & 0.0782 & 0.0000 & 0.0137 & 0.0402 & 0.0937 & 0.7039\\
IMS\_D & 8,129 & 0.2292 & 0.4203 & 0 & 0 & 0 & 0 & 1\\
Size & 8,129 & 22.1675 & 1.1662 & 18.1225 & 21.3416 & 21.9554 & 22.8022 & 27.5351\\
Lev & 8,129 & 0.4049 & 0.2034 & 0.0162 & 0.2435 & 0.3914 & 0.5447 & 1.9391\\
ROA & 8,129 & 0.0218 & 0.0764 & -1.5834 & 0.0043 & 0.0279 & 0.0574 & 0.7586\\
CFO & 8,129 & 0.0450 & 0.0714 & -0.9965 & 0.0088 & 0.0429 & 0.0812 & 0.8385\\
Growth & 8,129 & 0.1103 & 0.7424 & -1.4449 & -0.0815 & 0.0523 & 0.2008 & 33.7671\\
FirmAgeLog & 8,129 & 2.1655 & 0.8124 & 0.0000 & 1.6094 & 2.3979 & 2.7081 & 3.4965\\
SOE & 8,129 & 0.2266 & 0.4187 & 0 & 0 & 0 & 0 & 1\\
Dual & 8,129 & 0.3635 & 0.4810 & 0 & 0 & 0 & 1 & 1\\
BoardSize & 8,129 & 2.0660 & 0.1968 & 1.3863 & 1.9459 & 2.1972 & 2.1972 & 2.7726\\
IndRatio & 8,129 & 0.3822 & 0.0543 & 0.2222 & 0.3333 & 0.3750 & 0.4286 & 0.7143\\
Big4 & 8,129 & 0.0432 & 0.2033 & 0 & 0 & 0 & 0 & 1\\
LnAuditFee & 8,129 & 13.9500 & 0.5622 & 12.2549 & 13.5798 & 13.8643 & 14.2786 & 17.1269\\
AuditorTenure & 8,129 & 7.7276 & 5.8769 & 1 & 3 & 6 & 11 & 34\\
IsChangeFirm & 8,129 & 0.1172 & 0.3217 & 0 & 0 & 0 & 0 & 1\\
\bottomrule
\end{longtable}
\end{landscape}

需要注意，Growth最大值为33.7671，Lev最大值为1.9391，ROA和CFO也存在较大的负向极端值。当前阶段尚未进行1\%和99\%缩尾，因此这些极端值可能影响OLS系数和标准误；这也是后续稳健性检验与当前目标之间的重要差距。

\section{基准回归列次设计及其目的}

回归表分为IMS与IMS\_D两个面板，每个面板各包含六列，并全部使用同一8,129个观测。

\begin{enumerate}
  \item \textbf{第（1）列：一元回归。}仅加入行业专长指标和常数项，不加入控制变量与固定效应。目的在于给出行业专长与漂绿之间的原始样本内关系，并与前期Pearson相关和分组检验相互印证。该列不能排除公司特征造成的混杂影响。
  \item \textbf{第（2）列：加入财务控制。}加入Size、Lev、ROA、CFO、Growth和FirmAgeLog。目的在于控制公司规模、资本结构、盈利能力、现金流、成长性和生命周期差异，检验原始关系是否只是公司财务特征所致。
  \item \textbf{第（3）列：进一步加入治理控制。}在第（2）列基础上加入SOE、Dual、BoardSize和IndRatio。目的在于控制产权性质与内部治理结构对漂绿行为及审计师选择可能产生的共同影响。
  \item \textbf{第（4）列：加入全部控制变量。}进一步加入Big4、LnAuditFee、AuditorTenure和IsChangeFirm。目的在于区分行业专长与事务所规模、审计投入、审计关系持续时间及事务所变更等审计特征的关系。
  \item \textbf{第（5）列：加入年度与行业固定效应。}在全部控制变量基础上加入Year FE和Industry FE，并按公司聚类标准误。这是本阶段的主模型，用于控制各年度共同冲击和行业平均差异，同时允许同一公司跨年度误差相关。
  \item \textbf{第（6）列：HC3标准误对照。}模型解释变量、控制变量和固定效应与第（5）列完全相同，仅将公司聚类标准误改为HC3异方差稳健标准误。该列用于观察统计显著性是否依赖标准误估计方式，不是一个新的经济模型；考虑面板结构，第（5）列仍应作为主要判断依据。
\end{enumerate}

上述六列先以IMS为核心解释变量重复估计，再以IMS\_D替换IMS。连续指标保留市场份额的全部信息，虚拟指标则检验达到10\%行业专家门槛的事务所是否与较低漂绿程度相关。二者结果方向一致，可避免结论仅由某一种度量方式驱动。

\section{基准回归结果}

表\ref{tab:reg}报告核心解释变量的逐步回归结果。完整控制变量系数已保存在同目录的回归结果文件中，本表聚焦阶段性研究问题及列次比较。

\begin{table}[htbp]
\centering
\caption{审计师行业专长与企业漂绿：共同样本基准回归}\label{tab:reg}
\resizebox{\textwidth}{!}{%
\begin{threeparttable}
\begin{tabular}{lcccccc}
\toprule
 & (1) & (2) & (3) & (4) & (5) & (6)\\
\midrule
\multicolumn{7}{l}{\textit{Panel A：连续行业专长指标IMS}}\\
IMS & $-1.0238^{***}$ & $-0.7200^{**}$ & $-0.7305^{**}$ & $-0.7477^{**}$ & $-0.8739^{***}$ & $-0.8739^{***}$\\
 & (0.3157) & (0.3170) & (0.3159) & (0.3175) & (0.3150) & (0.2200)\\
调整$R^2$ & 0.0035 & 0.0127 & 0.0131 & 0.0183 & 0.0187 & 0.0187\\
\midrule
\multicolumn{7}{l}{\textit{Panel B：行业专家虚拟变量IMS\_D}}\\
IMS\_D & $-0.1471^{***}$ & $-0.0992^{**}$ & $-0.1003^{**}$ & $-0.1030^{**}$ & $-0.1129^{**}$ & $-0.1129^{***}$\\
 & (0.0511) & (0.0506) & (0.0506) & (0.0506) & (0.0523) & (0.0383)\\
调整$R^2$ & 0.0020 & 0.0119 & 0.0123 & 0.0176 & 0.0177 & 0.0177\\
\midrule
财务控制 & 否 & 是 & 是 & 是 & 是 & 是\\
治理控制 & 否 & 否 & 是 & 是 & 是 & 是\\
审计控制 & 否 & 否 & 否 & 是 & 是 & 是\\
年度固定效应 & 否 & 否 & 否 & 否 & 是 & 是\\
行业固定效应 & 否 & 否 & 否 & 否 & 是 & 是\\
标准误 & 公司聚类 & 公司聚类 & 公司聚类 & 公司聚类 & 公司聚类 & HC3\\
观测值 & 8,129 & 8,129 & 8,129 & 8,129 & 8,129 & 8,129\\
\bottomrule
\end{tabular}
\begin{tablenotes}
\footnotesize
\item 注：括号内为标准误。第（1）--（5）列按公司聚类，第（6）列为HC3异方差稳健标准误。$^{***}$、$^{**}$、$^{*}$分别表示在1\%、5\%和10\%水平显著。各列均使用同一共同样本。
\end{tablenotes}
\end{threeparttable}}
\end{table}

\subsection{IMS结果}

第（1）列中，IMS系数为$-1.0238$，在1\%水平显著，说明未控制其他因素时，事务所行业市场份额越高，客户公司的漂绿程度越低。第（2）列加入财务控制后，系数绝对值降至$-0.7200$，但仍在5\%水平显著。这一变化说明公司财务特征解释了部分原始相关关系，但不能完全解释该关系。

第（3）列加入公司治理变量后，IMS系数为$-0.7305$；第（4）列加入审计特征后为$-0.7477$，均在5\%水平显著。治理和审计变量的逐步加入没有改变系数方向，且系数变化较小。第（5）列控制年度与行业固定效应后，IMS系数为$-0.8739$，公司聚类标准误为0.3150，$p=0.0055$，在1\%水平显著。若IMS提高0.1，即行业客户收入市场份额提高10个百分点，则连续漂绿指标平均降低约$0.0874$。第（6）列的系数与第（5）列相同，但HC3标准误降至0.2200，显著性仍为1\%。

IMS在所有列均显著为负，是当前结果中最稳定且最有意义的发现。尤其是固定共同样本后，从一元模型到完整模型的符号与显著性均未改变，表明结果并非由逐列样本变化造成。

\subsection{IMS\_D结果}

以IMS\_D替换IMS后，第（1）列系数为$-0.1471$，在1\%水平显著；加入财务控制后降为$-0.0992$，但仍在5\%水平显著。第（3）和第（4）列系数分别为$-0.1003$和$-0.1030$，均在5\%水平显著。

第（5）列中，IMS\_D系数为$-0.1129$，公司聚类标准误为0.0523，$p=0.0309$。这意味着在其他条件及年度、行业效应相同的情况下，由IMS达到10\%门槛的行业专家事务所审计，与连续漂绿指标平均降低约0.1129相关。第（6）列使用HC3标准误后，系数不变且在1\%水平显著。IMS\_D与IMS结果方向一致，使得负向关系不依赖于连续或虚拟行业专长口径。

\subsection{控制变量结果中值得关注之处}

在最终公司聚类模型中，公司规模系数显著为负，资产负债率系数为正且至少弱显著；审计任期系数显著为正，事务所变更系数显著为正。Big4和审计收费在公司聚类标准误下未达到常用显著性水平，但在HC3标准误下显著为负。ROA、CFO、SOE、Dual、BoardSize和IndRatio在当前最终模型中未表现出稳定显著关系。

审计任期和事务所变更均与更高漂绿程度显著相关，是当前结果中值得后续核查的控制变量表现。不过，本阶段没有围绕这些控制变量提出研究假设，也未进行专门识别，因而只能将其视为样本内条件相关，不能扩展为独立结论。Big4与审计收费的显著性依赖标准误类型，说明同一公司的组内相关会影响统计推断，也进一步支持正文以公司聚类结果为主。

\section{阶段性结论、问题与尚未完成的目标}

\subsection{现阶段能够得到的结论}

在统一8,129个完整观测、逐步加入财务、治理和审计控制变量，并进一步控制年度和行业固定效应后，IMS和IMS\_D均与企业连续漂绿指标显著负相关。连续IMS结果在各列至少达到5\%显著水平，固定效应主模型达到1\%；IMS\_D结果在各列至少达到5\%显著水平。现阶段证据与“审计师行业专长越高，企业漂绿程度越低”的研究预期一致。

\subsection{当前结果中存在的问题}

第一，模型调整$R^2$较低，最终模型约为0.018，说明现有解释变量仅解释了漂绿指标较小部分的样本变异。这不否定核心系数，但表明企业漂绿仍受到模型外因素影响，当前模型的整体解释力有限。

第二，连续变量尚未缩尾。Growth最大值为33.7671，Lev最大值超过1，ROA和CFO也存在极端值。因此，核心结果是否受少数异常观测影响尚未检验。

第三，公司聚类与HC3标准误下部分控制变量显著性不同。核心变量方向和显著性总体稳定，但Big4、审计收费等结果对推断方式敏感。考虑公司年度面板的组内相关，第（5）列应作为主要结果，第（6）列仅为对照。

第四，本阶段控制了年度和行业固定效应，但未控制公司固定效应。当前估计仍可能受到不随时间变化的公司特征影响；同时，事务所行业专长与客户选择之间可能存在自选择和反向关系。

第五，当前只使用连续漂绿指标进行OLS，没有使用二元Greenwashing进行Logit或Probit回归，也尚未用IPS和IPS\_D替换行业专长度量。因此，结论目前只适用于本阶段变量口径。

第六，国内十大事务所变量因2024年全部缺失，未纳入当前基准控制变量。现有审计控制包含Big4、审计收费、审计任期和事务所变更，但尚未控制国内十大身份。

\subsection{与完整研究目标尚有的差距}

当前已完成的是基准相关关系检验，而非完整因果分析。与既定研究目标相比，尚未完成：（1）对连续变量进行1\%和99\%缩尾后重新估计；（2）以IPS和IPS\_D作为替代行业专长指标；（3）以二元Greenwashing进行Logit回归并报告平均边际效应；（4）开展制造业与非制造业、分年度异质性检验；（5）进行VIF和其他模型诊断；（6）处理审计师选择的内生性与潜在反向因果。因此，当前报告应表述为“审计师行业专长与企业漂绿程度显著负相关”，不宜表述为“审计师行业专长导致或抑制企业漂绿”。

\section{阶段性总结}

本阶段仿照既有审计师行业专长研究的多元回归设计，以共同样本逐步加入公司财务、治理、审计特征及年度、行业固定效应。结果显示，事务所行业市场份额IMS和10\%门槛行业专家变量IMS\_D的系数均稳定为负。主模型中，IMS系数为$-0.8739$并在1\%水平显著，IMS\_D系数为$-0.1129$并在5\%水平显著。该结果为审计师行业专长与企业漂绿之间的负向关系提供了阶段性证据，但低调整$R^2$、未缩尾极端值、尚未使用替代指标及缺乏内生性处理，均要求后续结果保持审慎解释。

\begin{thebibliography}{99}
\bibitem{chen2013} 陈小林，王玉涛，陈运森，2013：《事务所规模、审计行业专长与知情交易概率》，《会计研究》第2期：69--77。
\bibitem{liu2010} 刘文军，米莉，傅倞轩，2010：《审计师行业专长与审计质量——来自财务舞弊公司的经验证据》，《审计研究》第1期：47--53。
\bibitem{wang2018} 王生年，宋媛媛，徐亚飞，2018：《审计师行业专长缓解了资产误定价吗？》，《审计研究》第2期：96--103。
\bibitem{li2021} 李姝，杜亚光，张晓哲，2021：《审计师行业专长与企业创新——基于管理层信息环境视角的分析》，《审计研究》第1期：106--115。
\bibitem{defond2000} DeFond, M. L., Francis, J. R., and Wong, T. J., 2000, ``Auditor Industry Specialization and Market Segmentation: Evidence from Hong Kong,'' \textit{Auditing: A Journal of Practice \& Theory}, 19(1): 49--66.
\bibitem{krishnan2003} Krishnan, G. V., 2003, ``Does Big 6 Auditor Industry Expertise Constrain Earnings Management?'' \textit{Accounting Horizons}, 17: 1--16.
\bibitem{reichelt2010} Reichelt, K. J., and Wang, D., 2010, ``National and Office-Specific Measures of Auditor Industry Expertise and Effects on Audit Quality,'' \textit{Journal of Accounting Research}, 48(3): 647--686.
\end{thebibliography}

\end{document}
